\documentclass[11pt,a4paper]{article}

% ---------- Paquetes ----------
\usepackage[utf8]{inputenc}
\usepackage[spanish,es-noindentfirst]{babel}
\usepackage[a4paper,margin=2.5cm,headheight=14pt]{geometry}
\usepackage{titlesec}
\usepackage{enumitem}
\usepackage{longtable}
\usepackage{array}
\usepackage{booktabs}
\usepackage{tabularx}
\usepackage{ragged2e}
\usepackage[table]{xcolor}
\usepackage{fancyhdr}
\usepackage{caption}
\usepackage{setspace}
\usepackage[hidelinks,colorlinks=false]{hyperref}
\usepackage[numbers,sort&compress]{natbib}
\def\UrlBreaks{\do\/\do\-\do.\do\%\do\?\do=\do\&\do\_\do\~}
\Urlmuskip=0mu plus 2mu

% ---------- Colores ----------
\definecolor{darkgreen}{RGB}{27,67,50}
\definecolor{midgreen}{RGB}{64,145,108}
\definecolor{lightgray}{RGB}{242,242,242}
\definecolor{riskhigh}{RGB}{247,202,201}
\definecolor{riskmed}{RGB}{255,230,180}
\definecolor{risklow}{RGB}{212,237,218}

% ---------- Encabezados ----------
\pagestyle{fancy}
\fancyhf{}
\renewcommand{\headrulewidth}{0.4pt}
\fancyhead[L]{\footnotesize Monitoreo Agrario --- Altiplano Cundiboyacense}
\fancyhead[R]{\footnotesize Entrega No.\ 1}
\fancyfoot[C]{\small\thepage}

% ---------- Secciones ----------
\titleformat{\section}{\normalfont\Large\bfseries\color{darkgreen}}{\thesection.}{0.6em}{}
\titleformat{\subsection}{\normalfont\large\bfseries\color{midgreen}}{\thesubsection}{0.6em}{}
\titleformat{\subsubsection}{\normalfont\bfseries}{\thesubsubsection}{0.6em}{}
\titlespacing*{\section}{0pt}{1.4em}{0.8em}

\newcommand{\supuesto}[1]{\textit{Supuesto del equipo: #1}}

\newcolumntype{Y}{>{\RaggedRight\arraybackslash}X}
\newcolumntype{L}[1]{>{\RaggedRight\arraybackslash}p{#1}}

\setlist[itemize]{topsep=2pt,itemsep=2pt,parsep=0pt}
\setstretch{1.08}

\begin{document}

% ======================================================
% PORTADA
% ======================================================
\begin{titlepage}
\centering
\vspace*{1.5cm}
{\Large\bfseries FORMULACIÓN Y EVALUACIÓN DE PROYECTOS\par}
\vspace{0.6cm}
{\large Entrega No.\ 1 --- Formulación del Proyecto y Evaluación de Viabilidad\par}
\vspace{2.2cm}

\rule{0.8\textwidth}{0.6pt}\\[0.5cm]
{\Huge\bfseries\color{darkgreen} Sistema de Monitoreo Agrario\par}
{\Huge\bfseries\color{darkgreen} para el Altiplano Cundiboyacense\par}
\vspace{0.3cm}
{\large\itshape Monitoreo agroclimático basado en sensores de papel biodegradables\par}
\rule{0.8\textwidth}{0.6pt}\\[1.5cm]

{\large\bfseries Integrantes del equipo\par}
\vspace{0.4cm}
{\large
Julian David Cabrera Barragán\\
Andrés Felipe Carvajal Forero\\
Angelo Alejandro Ibañez Hernández\\
Juan Felipe Chibuque López\\
Joan Francisco Gómez\par
}

\vfill
{\large Bogotá D.C., Colombia\\ Octubre de 2026\par}
\end{titlepage}

\tableofcontents
\clearpage

% ======================================================
\section{Resumen ejecutivo}
% ======================================================

El presente documento formula, como proyecto de inversión, el \textbf{Sistema de Monitoreo Agrario para el Altiplano Cundiboyacense}: una red de captura periódica de información agroclimática (humedad del suelo, temperatura y variables asociadas a eventos de helada) soportada en \textbf{sensores de papel biodegradables} de bajo costo, dirigida a pequeños y medianos productores agrícolas de clima frío en Cundinamarca y Boyacá.

La iniciativa responde a una problemática verificable: los productores del Altiplano Cundiboyacense --- en particular los dedicados al cultivo de papa y a la ganadería de leche --- enfrentan pérdidas recurrentes por variabilidad climática (heladas y sequías asociadas a El Niño/La Niña), al tiempo que carecen de información agroclimática oportuna y de bajo costo a escala de predio, en un contexto de brecha digital rural todavía amplia \citep{boyaca7dias2026,semana2020,mintic2023}. Frente a esto, el sistema propone una alternativa diferenciada respecto de las soluciones de agricultura de precisión disponibles en el mercado (sensores electrónicos convencionales, estaciones meteorológicas, boletines agroclimáticos institucionales): sensores desechables, compostables y de fabricación simple, pensados para productores de pequeña escala con baja disposición a pagar por equipos electrónicos permanentes.

Este documento desarrolla la idea de proyecto, el planteamiento del problema, la justificación, el alcance, los objetivos, los actores y beneficiarios, la propuesta de valor, el contexto territorial, la alineación estratégica con los instrumentos de planeación nacionales, departamentales y municipales, las alternativas existentes, los estudios requeridos (mercado, técnico, organizacional, legal, ambiental y financiero), los requerimientos de información para la evaluación financiera, los indicadores de evaluación que se emplearán, los escenarios de sensibilidad previstos y una matriz preliminar de riesgos. Se incluyen además, como anexos, la ficha ejecutiva de formulación y el listado consolidado de supuestos.

Conforme al alcance definido por el curso, \textbf{esta entrega no desarrolla el sistema tecnológico ni presenta cifras financieras definitivas}; su resultado es una formulación estructurada que permite decidir, en las siguientes etapas, qué estudios adelantar para determinar la viabilidad comercial, técnica, organizacional, ambiental y financiera del proyecto.

% ======================================================
\section{Descripción de la idea}
% ======================================================

\subsection{¿Qué es el proyecto?}
Es un proyecto de inversión en etapa de formulación que propone diseñar e implementar un \textbf{sistema de monitoreo agrario de bajo costo}, compuesto por sensores de papel biodegradables desplegados en parcelas agrícolas, una red simple de recolección de datos (lectura periódica manual o semiautomática) y un mecanismo de consolidación y entrega de información a los productores y a terceros interesados (asociaciones, entidades territoriales).

\subsection{¿Qué solución propone?}
Sensores fabricados sobre sustratos de papel o celulosa --- impresos con tintas conductoras o indicadoras --- capaces de reaccionar a variables físicas del suelo y del ambiente (humedad, temperatura, presencia de condiciones de helada) y de degradarse de forma natural al cabo de su vida útil en campo, evitando así el residuo electrónico típico de los sensores IoT convencionales. Existe evidencia científica reciente que respalda la viabilidad técnica de esta familia de dispositivos: sensores capacitivos de humedad impresos sobre sustratos de papel biodegradable con nanofibrillas de celulosa han demostrado vida útil superior a 113 días en campo \citep{sciencedirect2025paper}, y distintos grupos de investigación han desarrollado sensores de celulosa para fertilidad, pH, humedad y actividad biológica del suelo \citep{nature2026cellulose,ieee2023biodeg,sciencedirect2024review,tudelft_sf,osaka_sf}.

\subsection{¿Qué problema u oportunidad aborda?}
Aborda la insuficiencia de información agroclimática oportuna y de bajo costo a escala de predio para productores pequeños y medianos del Altiplano Cundiboyacense, que hoy toman decisiones de siembra, riego y protección frente a heladas con información tardía, genérica o inexistente (ver Sección \ref{sec:problema}).

\subsection{¿A quién está dirigido?}
A pequeños y medianos productores agrícolas (principalmente de papa, hortalizas y ganadería de leche de clima frío), sus asociaciones y cooperativas, y de forma indirecta a entidades territoriales y gremios que requieren información agregada del territorio para la toma de decisiones (ver Sección \ref{sec:actores}).

\subsection{¿Dónde se implementaría?}
En el territorio definido como área objetivo dentro del Altiplano Cundiboyacense: provincias de Sabana Centro y Almeidas en Cundinamarca, y provincias Centro, Tundama y Sugamuxi en Boyacá (ver justificación territorial en la Sección \ref{sec:territorial}).

\subsection{¿Cuál es el papel de los sensores de papel biodegradables?}
Son el componente diferenciador del sistema: actúan como unidad mínima de captura de datos en campo. Su función es registrar, de forma periódica y descentralizada, variables de humedad y temperatura del suelo (y, por extensión, condiciones asociadas a riesgo de helada), para luego ser recolectados, leídos y agregados en información útil para el productor.

\subsection{¿Qué valor diferencial plantea la iniciativa?}
Frente a sensores electrónicos convencionales (plástico, metal, baterías de litio), el uso de sustratos de papel permite: (i) menor costo unitario de fabricación, (ii) biodegradabilidad y menor huella de residuo electrónico, y (iii) un modelo de reposición periódica (``sensor desechable'') en lugar de un activo electrónico permanente que exige mantenimiento, calibración y reemplazo de baterías --- barrera de adopción identificada en la literatura sobre agricultura de precisión en Colombia \citep{scielo2022precision}.

\subsection{¿Qué información se espera obtener mediante el sistema?}
Series periódicas de humedad del suelo y temperatura (incluyendo registro de eventos de temperatura cercana o inferior a 0 °C compatibles con heladas) a escala de parcela, agregables por vereda, municipio y provincia.

\subsection{¿Cómo podría utilizarse dicha información para apoyar la actividad agrícola?}
Para anticipar riesgo de helada y activar medidas de protección (riego por aspersión, coberturas, humos), para programar riego según humedad real del suelo, para alimentar los boletines agroclimáticos de asociaciones y entidades territoriales, y para sustentar decisiones de siembra y rotación de cultivos con información local en lugar de estimaciones genéricas.

% ======================================================
\section{Planteamiento del problema}
\label{sec:problema}
% ======================================================

\subsection{Contexto del sector agrícola}
El sector agropecuario colombiano continúa siendo una fuente relevante de empleo e ingresos rurales, pero enfrenta una exposición estructural a la variabilidad climática. La papa es uno de los cultivos transitorios de mayor relevancia nacional: en 2023 ocupó 143.696 hectáreas (24,1\,\% del área sembrada en tubérculos y plátano) con una producción de 2.931.116 toneladas (47,9\,\% de la producción de ese grupo) \citep{dane_ena2023}. Cundinamarca, Boyacá, Nariño y Antioquia concentran cerca del 90\,\% del área cultivada y de la producción nacional de papa \citep{agronet_papa}.

\subsection{Contexto específico del Altiplano Cundiboyacense}
El Altiplano Cundiboyacense es una altiplanicie de la cordillera Oriental, compartida por Cundinamarca y Boyacá, con una extensión aproximada de 25.000 km² y alturas entre 2.500 y 4.000 m s.\,n.\,m. \citep{wikipedia_altiplano}; comprende, entre otras, las subregiones de la Sabana de Bogotá, el valle de Ubaté-Chiquinquirá y el valle de Tunja-Duitama-Sogamoso. Es, a la vez, zona de alta producción de papa, hortalizas y leche, y zona de páramo y de recarga hídrica sujeta a protección ambiental especial \citep{corpoboyaca_paramo,corpoboyaca_car_paramo}. En Cundinamarca, la frontera agrícola departamental supera 1,4 millones de hectáreas (62,3\,\% del territorio), distribuida en 116 municipios, de los cuales 60 reportaron siembra de papa en 2022 \citep{fedepapa_cundinamarca2023}; Boyacá, por su parte, es el segundo departamento productor de papa del país, con cerca de 33.000 ha sembradas en promedio por año \citep{fedepapa_boyaca2023}.

\subsection{Situación actual}
Cundinamarca sembró 41.960 ha de papa en 2022 (con una contracción proyectada de -7{,}0\,\% para 2023) y alcanzó un rendimiento promedio de 24{,}16 t/ha \citep{fedepapa_cundinamarca2023}; Boyacá sembró 30.951 ha en 2022, con un crecimiento proyectado de apenas 0{,}2\,\% para 2023 \citep{fedepapa_boyaca2023}. Estos crecimientos marginales --- o contracciones --- ocurren en un departamento (Cundinamarca) que dispone de 473.675 ha aptas para papa, de las cuales solo el 52{,}1\,\% tiene la mejor aptitud productiva \citep{fedepapa_cundinamarca2023}, lo que sugiere que la restricción no es únicamente de tierra disponible sino de condiciones de producción (incluida la información climática oportuna).

\subsection{Necesidad identificada}
Existe una necesidad de información agroclimática oportuna y desagregada a escala de predio. El Boletín Agroclimático Nacional --- liderado por el Ministerio de Agricultura y Desarrollo Rural (MADR), el IDEAM y la FAO --- se construye con periodicidad mensual y escala subregional \citep{ideam_ban}, por lo que no sustituye la necesidad de un dato de campo, de bajo costo y en tiempo cercano al real, al nivel de la parcela individual. Esta brecha se agrava por el rezago de conectividad rural: según la Encuesta Nacional de Calidad de Vida 2023, si bien el 65{,}6\,\% de los hogares del país cuenta con conexión a internet, persiste una brecha superior a 20 puntos porcentuales entre cabeceras urbanas y zonas rurales, y entre los hogares sin conexión el 48{,}6\,\% señala el costo como barrera principal \citep{mintic2023}.

\subsection{Población afectada}
Pequeños y medianos productores de papa, hortalizas y leche de clima frío en los municipios del Altiplano Cundiboyacense, en particular aquellos sin acceso a estaciones meteorológicas propias, sensores electrónicos comerciales o personal técnico permanente.

\subsection{Consecuencias del problema}
Las heladas son el evento climático de mayor recurrencia documentada en la zona. En 2024, Boyacá reportó 94.111 hectáreas afectadas por heladas --- con declaratoria de calamidad pública --- frente a cifras superiores a 5.000 ha en 2023; el número de municipios con afectaciones críticas pasó de 24 en 2023 a 59 en 2024, y del total de hectáreas afectadas el 18\,\% corresponde a cultivos de papa (82\,\% a pastos de ganadería), concentrándose las mayores pérdidas en las provincias de Sugamuxi, Tundama y Centro \citep{boyaca7dias2026}. En un episodio previo de heladas severas, cerca de 30.000 productores resultaron afectados entre Boyacá (12.500) y Cundinamarca (17.900), según cifras reportadas por las gobernaciones departamentales \citep{semana2020}. El IDEAM ha emitido alertas en las que departamentos como Boyacá y Cundinamarca concentran un alto número de municipios en alerta naranja por heladas \citep{elpais_ideam}.

\subsection{Limitaciones de las alternativas actuales}
Las alternativas de monitoreo disponibles --- estaciones meteorológicas institucionales, boletines agroclimáticos de escala regional y sensores IoT comerciales (p.\,ej. VIT Sensors, PRISMAB) --- presentan limitaciones de escala (información subregional, no de predio), de costo (equipos electrónicos de inversión y mantenimiento permanente) o de cobertura territorial, factores identificados en la literatura como las principales barreras a la adopción de agricultura de precisión por parte de pequeños productores en Colombia: ausencia de sistemas de recolección y gestión de información, dificultad de adaptación tecnológica al agroecosistema local y altos costos de inversión inicial \citep{scielo2022precision}.

\subsection{Oportunidad que podría ser atendida mediante el proyecto}
Existe una oportunidad de ofrecer un dispositivo de captura de datos de muy bajo costo marginal (por ser desechable y de fabricación simple), ambientalmente compatible con una zona de páramo y recarga hídrica protegida, y dirigido específicamente al segmento de productores que hoy queda fuera del alcance de la agricultura de precisión convencional.

\medskip
\noindent\supuesto{} el planteamiento anterior describe la problemática y la oportunidad; no constituye, en sí mismo, una descripción de la solución tecnológica, la cual se desarrolla en la Sección 2 y se profundizará en el estudio técnico de etapas posteriores.

% ======================================================
\section{Justificación}
% ======================================================

La formulación de este proyecto como iniciativa de inversión se sustenta en seis dimensiones de relevancia:

\subsection{Relevancia económica}
La papa y la ganadería de leche de clima frío son actividades económicas centrales para miles de familias rurales del Altiplano Cundiboyacense; Cundinamarca y Boyacá concentran, junto con Nariño y Antioquia, cerca del 90\,\% de la producción nacional de papa \citep{agronet_papa}. Reducir pérdidas por heladas y mejorar la programación del riego tiene impacto directo sobre el ingreso de estos productores.

\subsection{Relevancia productiva}
El rendimiento de la papa en Cundinamarca (24{,}16 t/ha) convive con una amplia disponibilidad de tierra apta no utilizada a su mejor nivel (menos de la mitad de las 473.675 ha aptas corresponde a la mejor categoría de aptitud) \citep{fedepapa_cundinamarca2023}, lo que sugiere margen de mejora productiva asociado, entre otros factores, a mejor información para la toma de decisiones agronómicas.

\subsection{Relevancia tecnológica}
El proyecto se inscribe en una línea de investigación internacional activa --- sensores de celulosa y papel para monitoreo descentralizado en agricultura de precisión \citep{nature2026cellulose,ieee2023biodeg,sciencedirect2025paper,sciencedirect2024review} --- que aún no tiene antecedentes de aplicación identificados en el contexto colombiano, lo que representa una oportunidad de apropiación tecnológica diferenciadora.

\subsection{Relevancia ambiental}
A diferencia de los sensores electrónicos convencionales, un sensor de papel biodegradable reduce el residuo electrónico en una zona de páramo y recarga hídrica sujeta a protección especial por parte de la CAR y Corpoboyacá \citep{corpoboyaca_paramo,corpoboyaca_car_paramo}, donde el uso del suelo agrícola ya está normativamente restringido en los sectores de mayor sensibilidad.

\subsection{Relevancia social}
El proyecto se dirige preferentemente a pequeños y medianos productores, el segmento con menor acceso actual a tecnologías de agricultura de precisión y con mayor brecha de conectividad \citep{mintic2023,scielo2022precision}, contribuyendo potencialmente a cerrar una brecha de acceso a información entre grandes y pequeños productores.

\subsection{Relevancia territorial}
El proyecto se alinea con los instrumentos de planeación vigentes en ambos departamentos, que priorizan explícitamente la tecnificación agropecuaria y la innovación aplicada al territorio (ver Sección \ref{sec:alineacion}), y con la institucionalidad agroclimática nacional (IDEAM--MADR--FAO) \citep{ideam_ban}.

\medskip
\noindent En conjunto, estas seis dimensiones evidencian que el problema identificado --- pérdidas climáticas recurrentes combinadas con una brecha de información agroclimática de bajo costo a escala de predio --- representa una oportunidad concreta para formular y, posteriormente, evaluar la viabilidad de este proyecto de inversión.

% ======================================================
\section{Alcance}
% ======================================================

\subsection{Alcance de esta entrega}
Esta entrega cubre exclusivamente la formulación inicial del proyecto y la definición de los estudios necesarios para su evaluación. No incluye el desarrollo técnico del sistema, la fabricación de sensores, la implementación de un piloto, la construcción de una plataforma tecnológica, la operación comercial, la presentación de estados financieros reales ni la implementación física del proyecto.

\subsection{Alcance del proyecto (una vez formulado y evaluado)}

\subsubsection{Incluye}
\begin{itemize}
\item Diseño funcional del sistema de monitoreo (sensores de papel biodegradables + mecanismo de recolección y consolidación de datos).
\item Definición del modelo de servicio dirigido a productores individuales, asociaciones y cooperativas agrícolas.
\item Estudios de mercado, técnico, organizacional, legal, ambiental y financiero necesarios para decidir sobre la viabilidad del proyecto.
\item Definición de indicadores agroclimáticos mínimos a capturar (humedad de suelo, temperatura, proxy de riesgo de helada).
\end{itemize}

\subsubsection{No incluye}
\begin{itemize}
\item Fabricación o validación de laboratorio de los sensores.
\item Desarrollo de software, plataforma de datos o aplicativos móviles.
\item Operación comercial, distribución o venta real del servicio.
\item Intervención física en predios o en zonas de páramo protegidas.
\item Trámites de permisos ambientales o registros ante autoridades (ICA, CAR, Corpoboyacá), que corresponden a etapas posteriores.
\end{itemize}

\subsection{Área geográfica}
\label{sec:area-geografica}
El territorio objetivo se define como las \textbf{provincias de Sabana Centro y Almeidas en Cundinamarca} (municipios de referencia: Zipaquirá, Cogua, Tausa, Villapinzón y Chocontá) y las \textbf{provincias de Centro, Tundama y Sugamuxi en Boyacá} (municipios de referencia: Tunja, Samacá, Cómbita, Duitama, Tibasosa, Sogamoso y Firavitoba). La selección se sustenta en tres criterios verificables: (i) estos municipios concentran los principales núcleos de producción de papa de ambos departamentos --- Villapinzón y Zipaquirá figuran entre los cuatro municipios de mayor producción de Cundinamarca \citep{portalnews_papa}, y varios de los municipios boyacenses listados coinciden con el área del Complejo de Páramo Altiplano Cundiboyacense \citep{corpoboyaca_paramo}; (ii) las provincias de Sugamuxi, Tundama y Centro son señaladas como las de mayor pérdida histórica por heladas en Boyacá \citep{boyaca7dias2026}; y (iii) corresponden a zonas de predominio de pequeño y mediano productor, antes que de agroindustria de gran escala.

\supuesto{} la lista de municipios de referencia es ilustrativa y deberá precisarse y validarse durante el estudio de mercado, incluyendo la cuantificación exacta de productores potenciales por municipio.

\subsection{Población objetivo}
Pequeños y medianos productores agrícolas (papa, hortalizas) y ganaderos de leche de clima frío del área geográfica definida, así como sus asociaciones y cooperativas, en calidad de potenciales clientes o usuarios directos del sistema.

% ======================================================
\section{Objetivos}
% ======================================================

\subsection{Objetivo general}
Formular una propuesta de proyecto de inversión para un sistema de monitoreo agrario basado en sensores de papel biodegradables, definiendo su problemática, alcance, objetivos, propuesta de valor y alineación estratégica, como base para su posterior evaluación de viabilidad comercial, técnica, organizacional, ambiental y financiera.

\subsection{Objetivos específicos}
\begin{enumerate}
\item Caracterizar la problemática de información agroclimática a escala de predio que enfrentan los pequeños y medianos productores del Altiplano Cundiboyacense.
\item Definir la idea de proyecto y su propuesta de valor diferenciadora frente a las alternativas de monitoreo agrícola existentes.
\item Delimitar el alcance territorial y poblacional de la iniciativa, justificando la selección del área objetivo dentro del Altiplano Cundiboyacense.
\item Identificar los actores y beneficiarios relacionados con el proyecto y construir la matriz de actores correspondiente.
\item Analizar la alineación del proyecto con los instrumentos de planeación nacional, departamental y municipal relevantes para el territorio.
\item Identificar las principales alternativas tecnológicas y comerciales existentes para el monitoreo agroclimático, a fin de establecer el grado de diferenciación de la propuesta.
\item Determinar los estudios de mercado, técnico, organizacional, legal, ambiental y financiero requeridos para evaluar la viabilidad del proyecto.
\item Establecer los indicadores financieros y las variables que deberán investigarse en la etapa de evaluación financiera, junto con un primer mapa de riesgos del proyecto.
\end{enumerate}

% ======================================================
\section{Actores y beneficiarios}
\label{sec:actores}
% ======================================================

La Tabla \ref{tab:actores} presenta la matriz de actores identificados para el proyecto.

\begin{longtable}{L{2.6cm}L{2.6cm}L{2.6cm}L{3cm}L{2.4cm}}
\caption{Matriz de actores del proyecto}\label{tab:actores}\\
\toprule
\textbf{Actor} & \textbf{Rol} & \textbf{Relación con la iniciativa} & \textbf{Necesidad / interés} & \textbf{Nivel de participación}\\
\midrule
\endfirsthead
\toprule
\textbf{Actor} & \textbf{Rol} & \textbf{Relación con la iniciativa} & \textbf{Necesidad / interés} & \textbf{Nivel de participación}\\
\midrule
\endhead
Productores agrícolas de papa y hortalizas (pequeña y mediana escala) & Usuario final / cliente potencial & Instalan y usan los sensores en sus predios & Información oportuna de humedad y riesgo de helada a bajo costo & Alto\\
Ganaderos de leche de clima frío & Usuario final / cliente potencial & Usuarios de información sobre condiciones de pastos y heladas & Reducir pérdidas de pastos por heladas & Medio\\
Asociaciones y cooperativas de productores (p.\,ej. articuladas con Fedepapa) & Cliente institucional / canal & Agregan demanda y podrían contratar el servicio a nombre de sus asociados & Información agregada para asesoría técnica a sus afiliados & Alto\\
Empresas agrícolas de mediana escala & Cliente potencial & Adoptantes tempranos con mayor capacidad de pago & Optimizar riego y reducir pérdidas & Medio\\
Entidades territoriales (gobernaciones de Cundinamarca y Boyacá, alcaldías, secretarías de agricultura) & Aliado institucional / posible cofinanciador & Interesadas en información agroclimática territorial para sus políticas agropecuarias & Insumos para programas de extensión agropecuaria (PDEA) & Medio\\
Autoridades ambientales (CAR, Corpoboyacá) & Regulador & Vigilan el uso del suelo en zonas de páramo y recarga hídrica & Garantizar que el despliegue de sensores no afecte zonas protegidas & Bajo (regulador, no ejecutor)\\
Proveedores de materiales e insumos (papel/celulosa, tintas conductoras, componentes mínimos) & Proveedor & Suministran insumos para la fabricación de los sensores & Demanda estable de insumos & Medio\\
Instituciones académicas (universidades regionales, p.\,ej.\ UPTC, Universidad de Cundinamarca) & Aliado técnico & Posible apoyo en investigación y validación de prototipos & Proyectos de investigación aplicada & Medio\\
Entidades de investigación y planificación (UPRA, Agrosavia) & Aliado técnico / fuente de información & Proveen información de aptitud de uso del suelo y ordenamiento productivo & Información territorial de calidad para orientar la producción & Bajo\\
Potenciales inversionistas (fondos de innovación agro, programas de MinCiencias/ADR) & Financiador potencial & Podrían financiar etapas de validación y escalamiento & Retorno de la inversión y/o impacto en innovación agropecuaria & Medio\\
\bottomrule
\end{longtable}

% ======================================================
\section{Propuesta de valor}
% ======================================================

La propuesta de valor del proyecto se estructura de la siguiente manera:

\begin{center}
\begin{tabularx}{\textwidth}{L{2.6cm}L{3.1cm}L{2.6cm}L{3cm}L{3cm}}
\toprule
\textbf{Problema identificado} & \textbf{Solución propuesta} & \textbf{Usuario / cliente} & \textbf{Beneficio esperado} & \textbf{Elemento diferenciador}\\
\midrule
Falta de información agroclimática oportuna y de bajo costo a escala de predio, frente a pérdidas recurrentes por heladas y variabilidad climática. & Sistema de monitoreo con sensores de papel biodegradables desplegados en parcelas, con lectura y consolidación periódica de datos. & Pequeños y medianos productores de papa, hortalizas y leche; sus asociaciones; entidades territoriales. & Mejor anticipación de heladas y programación de riego; reducción potencial de pérdidas productivas; información territorial agregada para gremios y entidades. & Sensores desechables, compostables y de bajo costo unitario, frente a sensores electrónicos permanentes que exigen mantenimiento, calibración y mayor inversión inicial.\\
\bottomrule
\end{tabularx}
\end{center}

\subsection{¿Qué diferencia a los sensores de papel biodegradables de las alternativas existentes?}
Las alternativas de agricultura de precisión disponibles en Colombia (p.\,ej.\ VIT Sensors, PRISMAB) ofrecen sensores electrónicos de humedad, temperatura y conductividad de alta precisión, con reducciones de consumo de agua reportadas de hasta 30\,\% \citep{vitsensors,prismab}; sin embargo, requieren una inversión de capital por unidad, mantenimiento y gestión de baterías, lo que las hace menos accesibles para productores de pequeña escala. Los sensores de papel biodegradables, en cambio, se conciben como un insumo de reposición periódica --- similar a un consumible agrícola --- de menor costo marginal y sin residuo electrónico, apoyados en desarrollos de investigación que demuestran su viabilidad técnica (vida útil en campo superior a 113 días para sensores capacitivos impresos en sustrato de papel con nanofibrillas de celulosa) \citep{sciencedirect2025paper}. Esta combinación de menor costo de entrada y menor impacto ambiental constituye el elemento diferenciador central de la propuesta frente a la oferta actual.

% ======================================================
\section{Contexto territorial}
\label{sec:territorial}
% ======================================================

\subsection{El Altiplano Cundiboyacense}
El Altiplano Cundiboyacense es una altiplanicie de la cordillera Oriental de los Andes colombianos, compartida por los departamentos de Cundinamarca y Boyacá, con una extensión aproximada de 25.000 km² y una elevación promedio de 2.600 m s.\,n.\,m. (rango 2.500--4.000 m) \citep{wikipedia_altiplano}. Comprende tres subregiones planas principales: la Sabana de Bogotá, el valle de Ubaté-Chiquinquirá y el valle de Tunja-Duitama-Sogamoso (alto Chicamocha) \citep{wikipedia_altiplano}. Gran parte de su periferia corresponde a ecosistemas de páramo bajo manejo de las autoridades ambientales regionales \citep{corpoboyaca_paramo}.

\subsection{Vocación agropecuaria del territorio}
El territorio es la principal zona papera del país: Cundinamarca y Boyacá concentran --- junto con Nariño y Antioquia --- cerca del 90\,\% del área cultivada y la producción nacional de papa \citep{agronet_papa}. En Cundinamarca, la frontera agrícola departamental supera 1,4 millones de hectáreas (62,3\,\% del territorio departamental), de las cuales 473.675 ha se consideran aptas para papa \citep{fedepapa_cundinamarca2023}. El clima frío de altiplano también sostiene una ganadería de leche relevante, actividad que --- junto con la papa --- concentra la mayor parte de las pérdidas reportadas por heladas (82\,\% en pastos y 18\,\% en papa, según el balance de 2024 en Boyacá) \citep{boyaca7dias2026}.

\subsection{Exposición climática}
El territorio es recurrentemente afectado por heladas, fenómeno intensificado en episodios de El Niño. El IDEAM emite boletines agroclimáticos mensuales --- en alianza con el MADR y la FAO --- precisamente por la relevancia de esta amenaza para la planeación agronómica del país \citep{ideam_ban}, y ha declarado alertas naranja con un número elevado de municipios de Boyacá y Cundinamarca en riesgo \citep{elpais_ideam}.

\subsection{Condición ambiental e institucional}
El territorio objetivo colinda o se superpone parcialmente con zonas de páramo protegidas, de competencia de la CAR (Cundinamarca) y Corpoboyacá, bajo el marco del Sistema Nacional Ambiental (Ley 99 de 1993) \citep{ley99_1993,car_acuerdo16_1998}; estas autoridades han adoptado planes de manejo ambiental conjuntos para los páramos compartidos entre ambos departamentos y han tomado medidas preventivas frente a actividad agrícola en sectores sensibles \citep{corpoboyaca_car_paramo}. Esta condición es relevante tanto como restricción (zonas de exclusión agrícola) como oportunidad (argumento ambiental a favor de sensores biodegradables frente a electrónicos).

% ======================================================
\section{Alineación estratégica}
\label{sec:alineacion}
% ======================================================

El proyecto se analiza frente a los instrumentos de planeación nacional, departamental y municipal relevantes. La Tabla \ref{tab:alineacion} presenta la matriz de trazabilidad estratégica.

\begingroup
\setlength{\tabcolsep}{4pt}
\begin{longtable}{L{2.5cm}L{2.9cm}L{2.9cm}L{2.8cm}L{2.6cm}}
\caption{Matriz de trazabilidad estratégica}\label{tab:alineacion}\\
\toprule
\textbf{Nivel} & \textbf{Instrumento de planificación} & \textbf{Objetivo / estrategia} & \textbf{Relación con el proyecto} & \textbf{Contribución del proyecto}\\
\midrule
\endfirsthead
\toprule
\textbf{Nivel} & \textbf{Instrumento de planificación} & \textbf{Objetivo / estrategia} & \textbf{Relación con el proyecto} & \textbf{Contribución del proyecto}\\
\midrule
\endhead
Nacional & Plan Nacional de Desarrollo 2022--2026 ``Colombia Potencia Mundial de la Vida'' (Ley 2294 de 2023) \citep{ley2294_2023,sac_pnd2023} & Transformación ``Ordenamiento del territorio alrededor del agua'' y fomento de ciencia, tecnología e innovación para el agro, con enfoque en agricultura y ganadería regenerativa y resiliencia climática & El sistema aporta información para un uso del agua y del suelo más eficiente en una zona de recarga hídrica, alineado con la lógica de ordenamiento alrededor del agua & Genera datos de humedad de suelo que permiten un uso más eficiente del recurso hídrico en el riego agrícola\\
Nacional & Plan de Ordenamiento Productivo para la Cadena de la Papa en Colombia, UPRA/MADR \citep{upra_papa} & Orientar la localización de la producción de papa según aptitud del suelo y mejorar su competitividad y sostenibilidad & El sistema provee información de campo (humedad, temperatura) complementaria a los mapas de aptitud de uso del suelo que elabora UPRA & Aporta un insumo dinámico (dato en tiempo cercano al real) que complementa la información estática de aptitud de suelos\\
Departamental\newline(Cundinamarca) & Plan de Desarrollo Departamental 2024--2027 ``Gobernando: más que un plan'' \citep{cundinamarca_pdd2024} y Plan Departamental de Extensión Agropecuaria --- PDEA 2024--2028 \citep{adr_pdea_cundinamarca} & Línea estratégica de competitividad compartida con tecnificación agropecuaria y extensión rural & El sistema constituye una herramienta de base tecnológica que puede integrarse a los servicios de extensión agropecuaria departamental & Provee datos agroclimáticos de campo que pueden nutrir las actividades de asistencia técnica del PDEA\\
Departamental\newline(Boyacá) & Plan de Desarrollo Departamental 2024--2027 ``Nuestro Gran Plan es Boyacá'', eje ``Desarrollo Agropecuario como eje de la Boyacá Grande'' y dimensión de I+D+i aplicada al territorio \citep{boyaca_pdd2024} & Fortalecer la agroindustria y la innovación empresarial aplicada al territorio (p.\,ej.\ Red InnovAcción) & El proyecto es, en sí mismo, una innovación tecnológica aplicada al territorio agropecuario boyacense & Puede postularse como iniciativa piloto dentro de los programas departamentales de innovación agropecuaria\\
Municipal & Planes de Desarrollo Municipal y EOT/PBOT de los municipios del área objetivo (p.\,ej.\ Zipaquirá, Villapinzón, Tunja, Sogamoso) & Por definir según cada instrumento municipal vigente & Pendiente de verificación específica & \supuesto{} no se identificaron citas textuales verificables de los planes municipales durante esta formulación; su revisión puntual queda como tarea pendiente para la siguiente etapa\\
\bottomrule
\end{longtable}
\endgroup

\noindent La alineación no se limita a mencionar los instrumentos: en cada fila de la Tabla \ref{tab:alineacion} se explicita el mecanismo concreto mediante el cual el proyecto contribuiría a los objetivos o estrategias del instrumento correspondiente (dato de humedad para uso eficiente del agua, insumo dinámico para mapas de aptitud, integración a servicios de extensión agropecuaria departamental, y postulación como iniciativa de innovación territorial).

% ======================================================
\section{Alternativas existentes}
% ======================================================

Se identificaron las siguientes categorías de alternativas para atender la problemática de información agroclimática:

\begin{itemize}
\item \textbf{Boletines agroclimáticos institucionales}: el Boletín Agroclimático Nacional (IDEAM--MADR--FAO) ofrece información de periodicidad mensual y alcance subregional \citep{ideam_ban}; es gratuito pero no llega a la escala de predio individual.
\item \textbf{Sensores IoT comerciales}: empresas como VIT Sensors y PRISMAB ofrecen sensores electrónicos de humedad, temperatura y conductividad, con reportes de ahorro de agua de hasta 30\,\% \citep{vitsensors,prismab}; representan una solución de mayor precisión pero de mayor costo de inversión y mantenimiento por unidad.
\item \textbf{Desarrollos académicos locales}: prototipos universitarios de monitoreo con redes LoRa y control remoto de riego, documentados en proyectos de la Universidad Distrital y la Universidad Cooperativa de Colombia, entre otros, aplicados principalmente a cultivos de mayor escala (p.\,ej.\ palma) \citep{redalyc_iot}.
\item \textbf{Métodos tradicionales}: observación empírica del productor, recomendaciones de vecinos o de asistencia técnica esporádica, y medidas reactivas frente a heladas (riego de emergencia, quema de materiales) sin anticipación basada en datos.
\item \textbf{Sensores biodegradables en fase de investigación (internacional)}: desarrollos recientes de sensores de celulosa y papel para monitoreo agrícola descentralizado, incluyendo variantes alimentadas por baterías de papel o por microorganismos del suelo, documentados en literatura científica reciente pero sin reporte identificado de aplicación comercial en Colombia \citep{nature2026cellulose,ieee2023biodeg,sciencedirect2025paper,sciencedirect2024review,tudelft_sf,osaka_sf,natureworldnews2023}.
\end{itemize}

\noindent La revisión sugiere que la propuesta sí representa una alternativa diferenciada: combina la escala de predio (que no ofrecen los boletines institucionales) con un costo de entrada significativamente menor al de los sensores IoT comerciales, apoyándose en una línea de investigación internacional activa que, hasta donde pudo verificarse en esta formulación, no registra antecedentes de aplicación en el contexto colombiano.

% ======================================================
\section{Estudios requeridos para la formulación del proyecto}
% ======================================================

\subsection{Estudio de mercado}
Deberá determinar: tamaño del mercado objetivo (número de UPA --- unidades productoras agropecuarias --- de papa, hortalizas y leche en el área objetivo); demanda potencial y disposición a pagar de productores individuales y asociaciones; oferta y precios de alternativas sustitutas (sensores IoT, asesoría técnica tradicional); canales de comercialización viables (venta directa, a través de asociaciones/cooperativas, o mediante convenios con entidades territoriales); y potencial de adopción, considerando las barreras de brecha digital y de acceso a tecnología identificadas en la Sección \ref{sec:problema} \citep{mintic2023,scielo2022precision}.

\subsection{Estudio técnico}
Deberá investigar: tecnología y materiales de fabricación de los sensores de papel biodegradables (sustratos de celulosa, tintas conductoras, mecanismos de transducción) con base en la literatura técnica disponible \citep{nature2026cellulose,ieee2023biodeg,sciencedirect2025paper}; procesos de producción a escala piloto; infraestructura mínima requerida (producción, almacenamiento, logística de distribución y recolección en campo); capacidad de producción necesaria según el tamaño de mercado estimado; recursos humanos y técnicos necesarios; y esquema de operación y mantenimiento del sistema (incluida la vida útil esperada del sensor en campo, reportada en la literatura como superior a 113 días para formulaciones similares) \citep{sciencedirect2025paper}.

\subsection{Estudio organizacional}
Deberá determinar la estructura organizativa, el personal y las capacidades necesarias para fabricar, distribuir, operar y dar soporte al sistema, incluyendo si la producción de sensores se realizaría de forma directa o mediante terceros, y qué estructura se requeriría para la relación comercial con asociaciones y entidades territoriales.

\subsection{Estudio legal}
Deberá identificar, como mínimo: normativa ambiental aplicable al uso de suelo agrícola en zonas de páramo y de recarga hídrica, de competencia de la CAR y Corpoboyacá, bajo el marco de la Ley 99 de 1993 \citep{ley99_1993,car_acuerdo16_1998,corpoboyaca_car_paramo}; régimen de protección de datos personales de los productores que se vincule al servicio (Ley 1581 de 2012) \citep{ley1581_2012}, en caso de que el sistema recolecte datos asociados a predios o personas identificables; requisitos de registro o permisos ante el ICA si llegaran a emplearse insumos regulados; y normas de ordenamiento del uso del suelo agropecuario (UPRA) aplicables al área objetivo \citep{upra_papa}.

\subsection{Estudio ambiental}
Deberá evaluar el impacto ambiental del despliegue y retiro de los sensores en campo, y en particular el aporte neto del uso de sensores biodegradables (menor generación de residuo electrónico) frente a la alternativa de sensores electrónicos convencionales, considerando la condición de páramo y recarga hídrica de buena parte del territorio objetivo \citep{corpoboyaca_paramo,corpoboyaca_car_paramo}. Deberá documentarse igualmente el proceso de biodegradación esperado de los materiales (tiempos, condiciones de suelo) con base en la evidencia técnica disponible \citep{sciencedirect2024review,osaka_sf}.

\subsection{Estudio financiero}
Deberá definir la información necesaria para evaluar inversión inicial, costos (producción, operación, administración), ingresos esperados según los distintos modelos de comercialización identificados en el estudio de mercado, fuentes de financiamiento, flujo de caja, rentabilidad, punto de equilibrio y riesgos financieros (ver desarrollo en la Sección \ref{sec:financiero}). El marco general de este estudio deberá ser consistente con la lógica de formulación y evaluación de proyectos de inversión empleada en Colombia, tanto en su vertiente de evaluación privada \citep{sapag2011} como en la metodología de referencia para proyectos de inversión pública del DNP (MGA), en caso de explorarse fuentes de cofinanciación pública \citep{dnp_mga2023}.

% ======================================================
\section{Requerimientos para la evaluación financiera}
\label{sec:financiero}
% ======================================================

La Tabla \ref{tab:financiero} resume las variables que deberán investigarse y construirse en la etapa de evaluación financiera. Consistente con el alcance de esta entrega, \textbf{no se presentan valores definitivos}, sino la información requerida y su uso previsto.

\begin{longtable}{L{3.2cm}L{5.6cm}L{5.6cm}}
\caption{Requerimientos de información para la evaluación financiera}\label{tab:financiero}\\
\toprule
\textbf{Variable} & \textbf{Información requerida} & \textbf{Uso previsto}\\
\midrule
\endfirsthead
\toprule
\textbf{Variable} & \textbf{Información requerida} & \textbf{Uso previsto}\\
\midrule
\endhead
Horizonte de evaluación & Vida útil estimada del proyecto y periodo de análisis financiero. \supuesto{} se adoptará preliminarmente un horizonte de 10 años, propio de proyectos de inversión de infraestructura/servicio de baja obsolescencia, sujeto a validación técnica. & Definir el número de periodos del flujo de caja proyectado.\\
Inversión inicial & Costos de diseño de los sensores, capacidad de producción inicial, equipos mínimos de recolección/lectura de datos y capital de arranque. & Base del cálculo de la inversión (CAPEX) para el flujo de caja del periodo 0.\\
Costos de operación & Costos de distribución, recolección, logística en campo y soporte a productores. & Proyección de egresos operativos recurrentes.\\
Costos de producción & Costo unitario del sensor de papel biodegradable (materiales, tintas, mano de obra), a determinar en el estudio técnico. & Cálculo del costo variable por unidad producida/vendida.\\
Costos administrativos & Estructura mínima de personal administrativo y de gestión comercial. & Proyección de gastos fijos administrativos.\\
Ingresos esperados & Precio por sensor o por servicio (suscripción), volumen esperado de productores/asociaciones vinculados, según el estudio de mercado. & Proyección de ingresos del flujo de caja.\\
Capital de trabajo & Recursos necesarios para sostener el ciclo de producción-distribución-cobro. & Ajuste del flujo de caja por necesidades de liquidez operativa.\\
Fuentes de financiación & Recursos propios de los socios, posible cofinanciación de programas de innovación agropecuaria (ADR, MinCiencias) o de la banca de desarrollo. & Definición de la estructura de capital y del costo de financiación.\\
Tasa de descuento & Costo de oportunidad del capital o costo promedio ponderado de capital (WACC). \supuesto{} se usará como referencia preliminar una tasa real entre 12\,\% y 14\,\% E.\,A., propia de proyectos agroindustriales de riesgo moderado-alto en Colombia, sujeta a estimación formal en la etapa financiera. & Descuento de los flujos de caja proyectados para el cálculo del VPN.\\
Impuestos & Régimen tributario aplicable según la figura jurídica que adopte el proyecto. & Ajuste de los flujos de caja después de impuestos.\\
Valor residual & Valor de los activos (equipos de producción, infraestructura) al final del horizonte de evaluación, si aplica. & Incorporación al flujo de caja del último periodo.\\
Flujo de caja & Consolidación periódica de ingresos y egresos proyectados. & Insumo central para el cálculo de todos los indicadores de evaluación (Sección \ref{sec:indicadores}).\\
\bottomrule
\end{longtable}

% ======================================================
\section{Indicadores de evaluación}
\label{sec:indicadores}
% ======================================================

\begin{longtable}{L{3.6cm}L{5.5cm}L{5.5cm}}
\caption{Indicadores de evaluación financiera}\label{tab:indicadores}\\
\toprule
\textbf{Indicador} & \textbf{Qué mide} & \textbf{Para qué será utilizado / cómo contribuye a la decisión}\\
\midrule
\endfirsthead
\toprule
\textbf{Indicador} & \textbf{Qué mide} & \textbf{Para qué será utilizado / cómo contribuye a la decisión}\\
\midrule
\endhead
Valor Presente Neto (VPN/VAN) & El valor actual, descontado a la tasa de oportunidad del capital, de los flujos de caja netos generados por el proyecto durante el horizonte de evaluación. & Determinar si el proyecto genera valor económico adicional al costo de oportunidad del capital invertido; un VPN positivo sustentará la recomendación de inversión.\\
Tasa Interna de Retorno (TIR) & La tasa de descuento que hace que el VPN del proyecto sea igual a cero. & Comparar la rentabilidad esperada del proyecto frente a la tasa de descuento mínima aceptable (costo de capital) definida por los inversionistas.\\
Período de recuperación de la inversión (Payback) & El tiempo requerido para que los flujos de caja acumulados igualen la inversión inicial. & Evaluar el riesgo de liquidez y el horizonte en el que los inversionistas recuperan el capital invertido.\\
Punto de equilibrio & El volumen de ventas (sensores/servicios) o el nivel de ingresos en el que los ingresos totales igualan los costos totales. & Identificar el volumen mínimo de operación necesario para no generar pérdidas, insumo clave dado el modelo de ingreso por consumible desechable.\\
Relación beneficio/costo (B/C) & La razón entre el valor presente de los beneficios y el valor presente de los costos del proyecto. & Complementar el VPN con una medida relativa de eficiencia, especialmente útil si se exploran fuentes de cofinanciación pública evaluadas bajo metodologías tipo MGA \citep{dnp_mga2023}.\\
\bottomrule
\end{longtable}

\noindent El marco metodológico general para el cálculo e interpretación de estos indicadores seguirá la literatura estándar de formulación y evaluación de proyectos de inversión \citep{sapag2011}.

% ======================================================
\section{Escenarios y sensibilidad}
% ======================================================

\subsection{Escenarios}
Se definirán, como mínimo, tres escenarios para la evaluación financiera posterior:
\begin{itemize}
\item \textbf{Escenario pesimista}: baja tasa de adopción por parte de productores, precios de venta ajustados a la baja por sensibilidad al costo, y persistencia de la brecha de conectividad/digitalización rural como barrera de adopción \citep{mintic2023}.
\item \textbf{Escenario base}: adopción gradual alineada con la tendencia de tecnificación agropecuaria impulsada por los planes departamentales de extensión agropecuaria \citep{adr_pdea_cundinamarca,boyaca_pdd2024}, y costos de producción conforme a lo documentado en la literatura técnica de referencia.
\item \textbf{Escenario optimista}: adopción acelerada por eventos de helada recurrentes que incrementen la disposición a pagar por información anticipada (en línea con los antecedentes de afectación masiva documentados) \citep{boyaca7dias2026,semana2020}, y posible cofinanciación institucional (ADR, MinCiencias, gobernaciones).
\end{itemize}

\subsection{Variables críticas de sensibilidad}
Se identifican como variables de mayor impacto potencial sobre la viabilidad financiera: precio del sensor o del servicio; demanda (número de productores/asociaciones vinculados); número de clientes institucionales (asociaciones, entidades territoriales); costos de producción del sensor; inversión inicial requerida; costos operativos de distribución y recolección en campo; y tasa de adopción tecnológica, dado que esta última ha sido identificada en la literatura como la principal barrera a superar para la agricultura de precisión en productores pequeños en Colombia \citep{scielo2022precision}.

% ======================================================
\section{Riesgos iniciales}
% ======================================================

La Tabla \ref{tab:riesgos} presenta la matriz de riesgos preliminar del proyecto.

\begin{longtable}{L{2.6cm}L{2.1cm}L{1.6cm}L{1.6cm}L{1.5cm}L{3.8cm}}
\caption{Matriz de riesgos preliminar}\label{tab:riesgos}\\
\toprule
\textbf{Riesgo} & \textbf{Categoría} & \textbf{Probab.} & \textbf{Impacto} & \textbf{Nivel} & \textbf{Tratamiento propuesto}\\
\midrule
\endfirsthead
\toprule
\textbf{Riesgo} & \textbf{Categoría} & \textbf{Probab.} & \textbf{Impacto} & \textbf{Nivel} & \textbf{Tratamiento propuesto}\\
\midrule
\endhead
Baja disposición a pagar de pequeños productores por un servicio nuevo & Comercial & Alta & Alto & \cellcolor{riskhigh}Alto & Validar disposición a pagar mediante estudio de mercado directo y explorar modelos de pago agregado vía asociaciones\\
Baja tasa de adopción tecnológica por brecha digital/rural & De adopción tecnológica & Alta & Alto & \cellcolor{riskhigh}Alto & Diseñar el sistema con requerimientos mínimos de conectividad y apoyarse en asociaciones como intermediarias de adopción\\
Vida útil del sensor insuficiente en condiciones reales de campo (humedad extrema, suelos ácidos) & Técnico & Media & Alto & \cellcolor{riskmed}Medio-Alto & Validar en estudio técnico con pruebas piloto controladas antes de escalar producción\\
Costo de producción superior al estimado, afectando el punto de equilibrio & Financiero & Media & Alto & \cellcolor{riskmed}Medio-Alto & Cotización temprana de insumos y diseño modular que permita ajustar el costo unitario\\
Restricciones normativas para actividad agrícola/instalación de dispositivos en zonas de páramo & Ambiental / Legal & Media & Alto & \cellcolor{riskmed}Medio-Alto & Delimitar el despliegue fuera de las zonas núcleo de páramo protegidas por CAR/Corpoboyacá y consultar previamente con estas autoridades \citep{corpoboyaca_paramo,corpoboyaca_car_paramo}\\
Incumplimiento del régimen de protección de datos personales de los productores & Legal & Baja & Medio & \cellcolor{risklow}Bajo-Medio & Diseñar el tratamiento de datos conforme a la Ley 1581 de 2012 desde la etapa de formulación técnica \citep{ley1581_2012}\\
Dificultad logística para la recolección periódica de sensores en zonas rurales dispersas & Operativo & Media & Medio & \cellcolor{riskmed}Medio & Diseñar rutas de recolección apoyadas en la estructura de las asociaciones de productores\\
Entrada de competidores (sensores IoT u otras soluciones) con mayor capacidad de inversión & Comercial & Media & Medio & \cellcolor{riskmed}Medio & Reforzar el posicionamiento de precio/sostenibilidad como ventaja diferencial frente a sensores electrónicos\\
Variabilidad climática extrema que supere la capacidad de anticipación del sistema (eventos no previstos por el modelo) & Ambiental & Media & Alto & \cellcolor{riskmed}Medio-Alto & Integrar las alertas del Boletín Agroclimático Nacional del IDEAM como capa complementaria de información \citep{ideam_ban}\\
Dependencia de pocos proveedores de insumos clave (sustratos de celulosa, tintas conductoras) & Operativo & Baja & Medio & \cellcolor{risklow}Bajo-Medio & Identificar y calificar más de un proveedor potencial durante el estudio técnico\\
Dificultad para obtener financiación en etapa temprana del proyecto & Financiero & Media & Alto & \cellcolor{riskmed}Medio-Alto & Explorar líneas de cofinanciación pública para innovación agropecuaria (ADR, MinCiencias) y estructurar el proyecto bajo metodologías reconocidas (MGA) \citep{dnp_mga2023}\\
\bottomrule
\end{longtable}

% ======================================================
\section{Conclusiones de la formulación}
% ======================================================

\begin{enumerate}
\item La problemática identificada --- pérdidas agrícolas recurrentes por heladas y variabilidad climática, combinadas con una brecha de información agroclimática a escala de predio y una brecha digital rural todavía amplia --- está respaldada por fuentes institucionales y gremiales verificables (DANE, Fedepapa, IDEAM, MinTIC), lo que sustenta la pertinencia de continuar con la evaluación del proyecto.
\item El Altiplano Cundiboyacense, y en particular las provincias seleccionadas como área objetivo, concentran una proporción significativa de la producción nacional de papa y de las pérdidas históricas por heladas, lo que confirma la relevancia territorial de la iniciativa.
\item Existe sustento científico internacional para la viabilidad técnica de sensores de papel/celulosa biodegradables aplicados al monitoreo agrícola, aunque no se identificaron antecedentes de aplicación comercial en Colombia, lo que representa tanto una oportunidad de diferenciación como un riesgo técnico a validar en el estudio técnico de la siguiente etapa.
\item El proyecto se alinea de forma concreta --- no meramente declarativa --- con instrumentos de planeación nacional y departamental vigentes en ambos departamentos, aunque la alineación a nivel municipal queda pendiente de verificación puntual.
\item Los estudios de mercado, técnico, legal y ambiental identificados en este documento son coherentes con la naturaleza del proyecto y constituyen la ruta crítica para determinar su viabilidad; en particular, el estudio técnico deberá validar la vida útil y el costo unitario real del sensor, variables que condicionan de forma directa el punto de equilibrio financiero.
\item Esta formulación establece qué información deberá obtenerse --- y cómo seguirá utilizándose --- para construir el modelo financiero del proyecto, sin pretender, en esta etapa, determinar su rentabilidad definitiva.
\end{enumerate}

\noindent En conjunto, el equipo considera que la formulación presentada responde de manera sustentada a la pregunta guía de esta entrega: existe un problema verificable (pérdidas climáticas e información agroclimática insuficiente a escala de predio), dirigido a una población objetivo definida (pequeños y medianos productores del Altiplano Cundiboyacense), mediante una propuesta diferenciada (sensores de papel biodegradables), con sentido territorial (alineación con instrumentos de planeación vigentes) y con una ruta clara de estudios pendientes para determinar su viabilidad integral.

% ======================================================
\section{Referencias}
% ======================================================

\noindent\supuesto{} el enunciado de la entrega no especifica un sistema de citación particular; el equipo adoptó un estilo numérico (citas entre corchetes en el texto, p.\,ej.\ [1]), con la referencia completa correspondiente en el listado a continuación, en el mismo orden en que aparecen citadas.

\begin{thebibliography}{99}\setlength{\itemsep}{6pt}

\bibitem{agronet_papa} Agronet / Ministerio de Agricultura y Desarrollo Rural. (s.\,f.). \textit{Cundinamarca, Boyacá, Nariño y Antioquia representan 90\,\% de la producción de papa}. \url{https://www.agronet.gov.co/Noticias/Paginas/Cundinamarca,-Boyac%C3%A1,-Nari%C3%B1o-y-Antioquia,-representan-90-de-la-producci%C3%B3n-de-papa.aspx}

\bibitem{adr_pdea_cundinamarca} Agencia de Desarrollo Rural (ADR). (2024). \textit{Plan Departamental de Extensión Agropecuaria --- PDEA Cundinamarca 2024--2028}. \url{https://www.adr.gov.co/wp-content/uploads/2024/10/PDEA_Cundinamarca_2024.pdf}

\bibitem{boyaca7dias2026} Boyacá 7 Días. (2026). \textit{Las heladas no dan tregua: agricultores de Boyacá en alerta}. \url{https://boyaca7dias.com.co/2026/01/05/las-helada-no-dan-tregua-agricultores-de-boyaca-en-alerta/}

\bibitem{boyaca_pdd2024} Gobernación de Boyacá. (2024). \textit{Plan de Desarrollo Departamental 2024--2027 ``Nuestro Gran Plan es Boyacá''}. \url{https://www.boyaca.gov.co/dependencia/gobernacion/planeacion-y-ejecucion/politicas-lineamientos-manuales/planes-estrategicos-politicas-lineamientos-manuales/plan-de-desarrollo-departamental/}

\bibitem{car_acuerdo16_1998} Corporación Autónoma Regional de Cundinamarca (CAR). (1998). \textit{Acuerdo No.\ 16 de 1998}. \url{https://repositoriocdim.esap.edu.co/bitstreams/c97c8e52-eb7c-452d-bb32-cc053ecacc5a/download}

\bibitem{corpoboyaca_paramo} Corpoboyacá. (s.\,f.). \textit{Páramo Altiplano Cundiboyacense}. \url{https://www.corpoboyaca.gov.co/sirap/paramo-altiplano-cundiboyacense/}

\bibitem{corpoboyaca_car_paramo} Noticias Día a Día. (s.\,f.). \textit{Cundinamarca fortalece la protección de páramos: CAR actúa ante riesgo hídrico en Boyacá}. \url{https://noticiasdiaadia.com/cundinamarca/cundinamarca-paramos-car}

\bibitem{cundinamarca_pdd2024} Gobernación de Cundinamarca. (2024). \textit{Plan de Desarrollo Departamental 2024--2027 ``Gobernando: más que un plan''}. \url{https://www.cundinamarca.gov.co/rendicion-de-cuentas/centro-documental/plan-de-sesarrollo}

\bibitem{dane_ena2023} Departamento Administrativo Nacional de Estadística (DANE). (2025). \textit{Encuesta Nacional Agropecuaria (ENA) 2023 --- Boletín técnico}. \url{https://www.dane.gov.co/files/operaciones/ENA/bol-ENA-2023.pdf}

\bibitem{dnp_mga2023} Departamento Nacional de Planeación (DNP). (2023). \textit{Manual conceptual de la Metodología General Ajustada (MGA)}. \url{https://colaboracion.dnp.gov.co/CDT/MGA/Tutoriales%20de%20funcionamiento/Manual%20conceptual.pdf}

\bibitem{elpais_ideam} El País (Cali). (s.\,f.). \textit{IDEAM lanza alerta por posibles heladas en varias partes de Colombia}. Consultado en octubre de 2026. \url{https://www.elpais.com.co/colombia/ideam-lanza-alerta-por-posibles-heladas-en-varias-partes-de-colombia.html}

\bibitem{fedepapa_boyaca2023} Federación Colombiana de Productores de Papa (Fedepapa). (2024). \textit{Boletín Regional Boyacá, Vol.\ 7 --- 2023}. \url{https://fedepapa.com/home/wp-content/uploads/2024/10/Regional-Boyaca.pdf}

\bibitem{fedepapa_cundinamarca2023} Federación Colombiana de Productores de Papa (Fedepapa). (2024). \textit{Boletín Regional Cundinamarca, Vol.\ 7 --- 2023}. \url{https://fedepapa.com/home/wp-content/uploads/2024/10/Regional-Cundinamarca.pdf}

\bibitem{ideam_ban} Instituto de Hidrología, Meteorología y Estudios Ambientales (IDEAM), Ministerio de Agricultura y Desarrollo Rural (MADR) y FAO. (s.\,f.). \textit{Boletín Agroclimático Nacional --- Protocolo versión 1.0}. \url{https://www.ideam.gov.co/sites/default/files/mapa-de-procesos/m-gdi-h-pc018_protocolo_ban_v1.pdf}

\bibitem{ieee2023biodeg} IEEE. (2023). \textit{Fabrication and Characterization of Soil Moisture Sensors on a Biodegradable, Cellulose-Based Substrate}. \url{https://ieeexplore.ieee.org/document/10220095/}

\bibitem{ley99_1993} Congreso de la República de Colombia. (1993). \textit{Ley 99 de 1993, por la cual se crea el Ministerio del Medio Ambiente y el Sistema Nacional Ambiental (SINA)}. \url{https://www.funcionpublica.gov.co/eva/gestornormativo/norma.php?i=297}

\bibitem{ley1581_2012} Congreso de la República de Colombia. (2012). \textit{Ley 1581 de 2012, Régimen General de Protección de Datos Personales}. \url{https://www.funcionpublica.gov.co/eva/gestornormativo/norma.php?i=49981}

\bibitem{ley2294_2023} Congreso de la República de Colombia. (2023). \textit{Ley 2294 de 2023, Plan Nacional de Desarrollo 2022--2026 ``Colombia Potencia Mundial de la Vida''}.

\bibitem{mintic2023} Ministerio de Tecnologías de la Información y las Comunicaciones (MinTIC). (2023). \textit{Índice de Brecha Digital 2023}. Colombia TIC. \url{https://colombiatic.mintic.gov.co/679/articles-396961_recurso_1.pdf}

\bibitem{nature2026cellulose} Nature Communications. (2026). \textit{Cellulose-based sensors for decentralized monitoring in precision agriculture}. \url{https://www.nature.com/articles/s41467-026-70730-7}

\bibitem{natureworldnews2023} Nature World News. (2023). \textit{Soil Moisture Sensing Goes Green With Biodegradable Technology for Sustainable Smart Agriculture}. \url{https://www.natureworldnews.com/articles/58976/20231018/soil-moisture-sensing-goes-green-biodegradable-technology-sustainable-smart-agriculture.htm}

\bibitem{portalnews_papa} Portal News. (2022). \textit{Los 4 municipios que más producen papa en Cundinamarca}. \url{https://portalnews.co/cundinamarca/los-4-municipios-que-mas-producen-papa-en-cundinamarca}

\bibitem{redalyc_iot} Redalyc. (s.\,f.). \textit{Internet de las cosas aplicado a la agricultura: estado actual}. \url{https://www.redalyc.org/journal/6139/613964509009/html/}

\bibitem{sac_pnd2023} Sociedad de Agricultores de Colombia (SAC). (2023). \textit{Plan Nacional de Desarrollo 2022--2026: el menú para los sectores agropecuario y rural}. \url{https://sac.org.co/plan-nacional-de-desarrollo-2022-2026-el-menu-para-los-sectores-agropecuario-y-rural/}

\bibitem{sapag2011} Sapag Chain, N. (2011). \textit{Proyectos de inversión: formulación y evaluación} (2.\textsuperscript{a} ed.). Pearson Educación.

\bibitem{scielo2022precision} Scielo. (2022). \textit{Variables que influyen en la aplicación de la agricultura de precisión en Colombia: revisión de estudios}. \url{http://www.scielo.org.co/scielo.php?script=sci_arttext&pid=S0122-87062022000100015}

\bibitem{sciencedirect2025paper} ScienceDirect. (2025). \textit{Inductively coupled capacitive soil moisture sensors printed on a biodegradable substrate: Characterization and long-term testing}. \url{https://www.sciencedirect.com/science/article/pii/S2772375525002850}

\bibitem{sciencedirect2024review} ScienceDirect. (2024). \textit{Review of low-cost, off-grid, biodegradable in situ autonomous soil moisture sensing systems: Is there a perfect solution?}. \url{https://www.sciencedirect.com/science/article/pii/S016816992400680X}

\bibitem{semana2020} Revista Semana. (2020). \textit{Heladas e incendios se mantendrán hasta marzo en Colombia}. \url{https://www.semana.com/medio-ambiente/articulo/heladas-e-incendios-se-mantendran-hasta-marzo-en-colombia/48539/}

\bibitem{tudelft_sf} TU Delft Research Portal. (s.\,f.). \textit{Design of biodegradable soil moisture sensor powered by bacteria charged paper batteries}. \url{https://research.tudelft.nl/en/publications/design-of-biodegradable-soil-moisture-sensor-powered-by-bacteria-/}

\bibitem{osaka_sf} Lab Horizons. (2025). \textit{Eco-Friendly Paper Sensors Use Case in Precise Farm Monitoring} (sobre el desarrollo de la Universidad de Osaka). \url{https://labhorizons.co.uk/2025/01/eco-friendly-paper-sensors-use-case-in-precise-farm-monitoring/}

\bibitem{upra_papa} Unidad de Planificación Rural Agropecuaria (UPRA). (2024). \textit{Plan de Ordenamiento Productivo para la Cadena de la Papa en Colombia}. \url{https://upra.gov.co/sites/default/files/2024-09/POP_Cadena_Papa.pdf}

\bibitem{vitsensors} VIT Sensors. (s.\,f.). \textit{Sensores Agrícolas Inteligentes para Monitoreo de Cultivos}. \url{https://www.vitsensors.com/sensores/}

\bibitem{prismab} PRISMAB. (s.\,f.). \textit{Sensores para la agricultura de precisión}. \url{https://prismab.com/}

\bibitem{wikipedia_altiplano} Wikipedia contributors. (2026). \textit{Altiplano Cundiboyacense}. Wikipedia, The Free Encyclopedia. [Fuente de caracterización geográfica general, de carácter tertiario; se referencia de forma complementaria a las fuentes institucionales citadas]. \url{https://en.wikipedia.org/wiki/Altiplano_Cundiboyacense}

\end{thebibliography}

% ======================================================
\section{Anexos}
% ======================================================

\subsection*{Anexo A. Ficha de formulación del proyecto}

\begin{longtable}{L{3.6cm}L{10.4cm}}
\toprule
\textbf{Elemento} & \textbf{Descripción}\\
\midrule
\endfirsthead
\toprule
\textbf{Elemento} & \textbf{Descripción}\\
\midrule
\endhead
Nombre del proyecto & Sistema de Monitoreo Agrario para el Altiplano Cundiboyacense\\
Problema & Pérdidas agrícolas recurrentes por heladas y variabilidad climática, combinadas con una brecha de información agroclimática oportuna a escala de predio y una brecha digital rural persistente, que afectan a pequeños y medianos productores del Altiplano Cundiboyacense\\
Idea & Sistema de monitoreo agrario basado en sensores de papel biodegradables para la captura periódica de humedad del suelo y temperatura (incluyendo proxy de riesgo de helada) a escala de parcela\\
Objetivo general & Formular una propuesta de proyecto de inversión para un sistema de monitoreo agrario basado en sensores de papel biodegradables, definiendo su problemática, alcance, objetivos, propuesta de valor y alineación estratégica, como base para su posterior evaluación de viabilidad\\
Población objetivo & Pequeños y medianos productores de papa, hortalizas y leche de clima frío, y sus asociaciones/cooperativas, en el área objetivo definida\\
Área geográfica & Provincias de Sabana Centro y Almeidas (Cundinamarca) y provincias Centro, Tundama y Sugamuxi (Boyacá), dentro del Altiplano Cundiboyacense\\
Propuesta de valor & Sensores de papel biodegradables, de bajo costo y reposición periódica, que proveen información agroclimática a escala de predio, frente a boletines institucionales de escala subregional y sensores IoT comerciales de mayor costo e inversión permanente\\
Beneficiarios & Productores agrícolas, ganaderos de leche, asociaciones/cooperativas, entidades territoriales y gremios del sector agropecuario del área objetivo\\
Diferenciador & Biodegradabilidad, bajo costo unitario y modelo de consumible desechable, frente a sensores electrónicos permanentes\\
Duración estimada & \supuesto{} 10 años de horizonte de evaluación (sujeto a validación en el estudio financiero); la fase de formulación y estudios (mercado, técnico, legal, ambiental, financiero) se estima en 6 a 12 meses adicionales a esta entrega\\
Resultado esperado & Una formulación estructurada del proyecto que permita decidir, con base en los estudios identificados, si es viable avanzar hacia el diseño técnico, la validación piloto y la evaluación financiera definitiva del sistema\\
\bottomrule
\end{longtable}

\subsection*{Anexo B. Listado consolidado de supuestos del equipo}

\begin{enumerate}
\item La lista de municipios de referencia del área objetivo (Sección \ref{sec:area-geografica}) es ilustrativa y deberá precisarse y validarse cuantitativamente en el estudio de mercado.
\item El planteamiento del problema (Sección \ref{sec:problema}) describe la problemática y la oportunidad sin asumir aún la solución tecnológica como validada técnicamente; dicha validación corresponde al estudio técnico.
\item La fila ``Municipal'' de la matriz de trazabilidad estratégica (Tabla \ref{tab:alineacion}) queda pendiente de verificación puntual contra los Planes de Desarrollo Municipal y EOT/PBOT vigentes de cada municipio del área objetivo.
\item El horizonte de evaluación financiera (10 años) y la tasa de descuento preliminar (12\,\%--14\,\% E.\,A.) presentados en la Tabla \ref{tab:financiero} son valores de referencia preliminares, no definitivos, sujetos a estimación formal en la etapa de evaluación financiera.
\item La duración estimada de la fase de estudios posteriores a esta entrega (6 a 12 meses) es una estimación preliminar del equipo, sujeta a ajuste según la planeación detallada del proyecto.
\item No se identificaron antecedentes verificables de aplicación comercial de sensores de papel biodegradables en el contexto agrícola colombiano; esta afirmación se basa en la búsqueda documental realizada para esta entrega y no constituye un estudio de mercado exhaustivo.
\end{enumerate}

\end{document}
